\chapter{切空间与微分形式}
\label{chap:tangent-differential-forms}

流形的坐标取值于欧氏空间，但同一点附近可以选用不同坐标。切向量、张量和
微分形式的定义若要反映流形自身的几何，就不能依赖某一张坐标图，也不能依赖
流形在欧氏空间中的某个嵌入。本章从曲线速度与导子两种观点构造切空间，继而
引入切映射、向量场、张量场和微分形式。最后把
\chapterref{chap:differential-forms-euclidean}中的外微分移植到流形上。
这些构造都不需要度量；对微分形式进行积分时，还要加入定向，这将在下一章
处理。

\section{切向量的等价定义}

设 \(M\) 为光滑 \(n\) 维流形，\(p\in M\)。考虑所有定义在 \(0\) 的某个
邻域且满足 \(\gamma(0)=p\) 的光滑曲线。若在一张含 \(p\) 的坐标图
\((U,x)\) 中
\[
 (x\circ\gamma_1)'(0)=(x\circ\gamma_2)'(0),
\]
就称 \(\gamma_1,\gamma_2\) 在 \(p\) 具有相同速度。这个条件与坐标图的
选择无关。事实上，若 \(y\) 是另一组坐标，则链式法则给出
\[
 (y\circ\gamma_j)'(0)
 =D(y\circ x^{-1})(x(p))(x\circ\gamma_j)'(0).
\]
曲线关于“具有相同速度”的等价类称为\term{曲线切向量}。

曲线速度类具有自然的线性运算。固定一张坐标图，把两个速度类的和及数乘定义
为坐标速度向量在 \(\R^n\) 中的和及数乘；向量 \(w\in\R^n\) 总能由曲线
\(t\mapsto x^{-1}(x(p)+tw)\) 在充分小的参数区间内实现。更换坐标时，速度
向量乘以坐标变换在 \(p\) 处的可逆导数。这个变换是线性的，所以刚才的运算
与坐标选择无关，曲线速度类由此成为向量空间。

还可以完全借助流形上的函数描述切向量。记 \(C_p^\infty(M)\) 为 \(p\)
附近光滑函数的芽所成的实代数：两个局部光滑函数在 \(p\) 的某个邻域相等
时，视为同一个芽。

\begin{definition}
\label{def:tangent-derivation}
\(p\) 处的\term{导子}是一个实线性映射
\[
 D:C_p^\infty(M)\longrightarrow\R,
\]
它满足 Leibniz 规则
\[
 D(fg)=f(p)Dg+g(p)Df.
\]
所有导子组成的向量空间记为 \(T_pM\)，称为 \(M\) 在 \(p\) 处的
\term{切空间}。
\end{definition}

函数芽只保留 \(p\) 附近的信息，正与切向量的局部性质相符。也可以从全局
光滑函数出发定义导子。为核对两种定义等价，先注意全局导子只依赖函数在
\(p\) 附近的取值：若 \(f\) 在 \(p\) 的一个邻域为零，取截断函数
\(\chi\) 使 \(\chi(p)=1\) 且 \(\operatorname{supp}\chi\) 包含在该邻域内，
则 \(\chi f=0\)，Leibniz 规则给
\[
 0=D(\chi f)=Df.
\]
反过来，任一函数芽可先缩小代表元的定义域，再乘以上述截断函数并补零，得到
全局光滑函数；两个这样的延拓在 \(p\) 附近相同，故全局导子给它们相同的值。
因此全局函数导子与函数芽导子确实得到同一个切空间。

曲线 \(\gamma\) 给出导子
\[
 D_\gamma f
 :=\left.\frac{d}{dt}\right|_{t=0}f(\gamma(t)).
\]
链式法则表明 \(D_\gamma\) 只依赖 \(\gamma\) 的速度类，乘积求导法则则表明
它满足 Leibniz 规则。

\begin{theorem}[切向量的曲线表示与导子表示]
\label{thm:tangent-curves-derivations}
从曲线速度类到导子的映射
\[
 [\gamma]\longmapsto D_\gamma
\]
是线性双射。在局部坐标 \((x^1,\ldots,x^n)\) 中，每个 \(v\in T_pM\)
唯一写成
\[
 v=\sum_{i=1}^n v^i
 \left.\frac{\partial}{\partial x^i}\right|_p,
 \qquad v^i=v(x^i).
\]
特别地，\(\dim T_pM=n\)。
\end{theorem}

\begin{proof}
先证明导子由它在坐标函数上的值决定。平移坐标后可设 \(x(p)=0\)，并把
坐标域缩小为以 \(0\) 为中心的开球。对任意局部光滑函数 \(f\)，记
\(\widetilde f=f\circ x^{-1}\)。微积分基本定理给出 Hadamard 分解
\[
 \widetilde f(u)-\widetilde f(0)
 =\sum_{i=1}^n u^i g_i(u),
 \qquad
 g_i(u)=\int_0^1\partial_i\widetilde f(tu)\,dt .
\]
各 \(g_i\) 光滑，并且
\(g_i(0)=\partial_i\widetilde f(0)\)。由
\(D1=D(1\cdot1)=2D1\) 可知导子在常数上为零。对 Hadamard 分解使用
Leibniz 规则，得到
\[
 Df=\sum_{i=1}^nD(x^i)\,\partial_i\widetilde f(0).
\]
因此导子的坐标表示存在且唯一。

反过来，给定 \(v=(v^1,\ldots,v^n)\in\R^n\)，在充分小的 \(t\) 上定义
\[
 \gamma_v(t)=x^{-1}(tv).
\]
链式法则给出
\[
 D_{\gamma_v}f
 =\sum_{i=1}^n v^i\partial_i\widetilde f(0),
\]
正好得到具有坐标分量 \(v^i\) 的导子。故每个导子都来自一条曲线，而两条
曲线给出同一导子，当且仅当它们的坐标速度相同。预先定义的速度类线性运算
在此对应于导子的逐点线性运算，所以上述双射是线性的。
\end{proof}

若 \(y=y(x)\) 是另一组坐标，则坐标基按照
\begin{equation}
 \left.\frac{\partial}{\partial x^i}\right|_p
 =\sum_{j=1}^n
 \frac{\partial y^j}{\partial x^i}(p)
 \left.\frac{\partial}{\partial y^j}\right|_p
 \label{eq:tangent-coordinate-change}
\end{equation}
变换。相应地，切向量的分量乘以坐标变换的 Jacobian。基与分量的变换互相
补偿，切向量本身不随坐标改变。

在欧氏开集上，这一定义恢复熟悉的方向向量：
\[
 T_p\R^n\simeq\R^n,\qquad
 v(f)=\sum_i v^i\partial_i f(p).
\]
对单位球面则有
\[
 T_pS^{n-1}=\{v\in\R^n:p\cdot v=0\}.
\]
确实，若 \(\gamma(t)\in S^{n-1}\) 且 \(\gamma(0)=p\)，对
\(\norm{\gamma(t)}^2=1\) 求导便得 \(p\cdot\gamma'(0)=0\)。反过来，若
\(p\cdot v=0\)，曲线
\[
 \gamma(t)=\frac{p+tv}{\norm{p+tv}}
\]
在零点的速度为 \(v\)。这也说明嵌入给出的环境向量表示很方便，但切空间的
定义本身并不需要先选一个嵌入。

\section{切映射、向量场和 Lie 括号}

若 \(F:M\to N\) 光滑，定义 \(F\) 在 \(p\) 处的\term{切映射}
\[
 dF_p:T_pM\longrightarrow T_{F(p)}N,\qquad
 (dF_pv)(g)=v(g\circ F).
\]
右端满足 Leibniz 规则，因而确实是 \(F(p)\) 处的导子。在曲线表示中，
\[
 dF_p[\gamma]=[F\circ\gamma],
\]
在局部坐标中，\(dF_p\) 就是由 \(F\) 的 Jacobian 矩阵表示的线性映射。

\begin{proposition}[切映射的函子性]
\label{prop:tangent-map-functoriality}
若 \(F:M\to N\)、\(G:N\to P\) 光滑，则
\[
 d(G\circ F)_p=dG_{F(p)}\circ dF_p,
 \qquad d(\operatorname{id}_M)_p=I_{T_pM}.
\]
若 \(F\) 是微分同胚，则 \(dF_p\) 在每一点都是线性同构，而且
\[
 (dF_p)^{-1}=d(F^{-1})_{F(p)}.
\]
\end{proposition}

\begin{proof}
对任意 \(h\in C_{G(F(p))}^\infty(P)\)，有
\[
 [dG_{F(p)}(dF_pv)]h
 =(dF_pv)(h\circ G)
 =v(h\circ G\circ F),
\]
这就是 \(d(G\circ F)_pv\) 对 \(h\) 的作用。恒等映射的结论直接来自
定义。若 \(F\) 是微分同胚，把函子性用于
\(F^{-1}\circ F=\operatorname{id}_M\) 和
\(F\circ F^{-1}=\operatorname{id}_N\)，便得到最后一个公式。
\end{proof}

\term{向量场}是在每一点 \(p\) 选取 \(X(p)\in T_pM\) 的光滑选择。所谓
光滑，是指在任一坐标图中
\[
 X=\sum_{i=1}^nX^i\frac{\partial}{\partial x^i}
\]
的系数 \(X^i\) 都光滑。向量场作用于函数所得
\[
 Xf(p)=X(p)f
\]
仍是光滑函数，并满足乘积规则。

两个向量场作为微分算子的复合 \(XY\) 一般含有二阶导数，不再是向量场；
交换子中的二阶项却会消去。定义
\[
 [X,Y]f=X(Yf)-Y(Xf).
\]

\begin{proposition}
\label{prop:lie-bracket-properties}
Lie 括号 \([X,Y]\) 是向量场，并在局部坐标中满足
\[
 [X,Y]
 =\sum_{k=1}^n
 \left(\sum_{i=1}^n
  \bigl(X^i\partial_iY^k-Y^i\partial_iX^k\bigr)\right)
 \frac{\partial}{\partial x^k}.
\]
它是双线性的、反对称的，并满足 Jacobi 恒等式
\[
 [X,[Y,Z]]+[Y,[Z,X]]+[Z,[X,Y]]=0.
\]
\end{proposition}

\begin{proof}
在坐标中展开 \(X(Yf)\) 和 \(Y(Xf)\)。含
\(\partial_i\partial_j f\) 的项因混合偏导对称而相消，只留下所列的一阶
项，所以交换子确实是向量场。双线性和反对称性直接来自定义。把向量场视为
函数代数上的线性算子，Jacobi 恒等式归结为算子交换子的恒等式
\[
 [A,[B,C]]+[B,[C,A]]+[C,[A,B]]=0;
\]
展开后，每个三重算子复合恰出现一次正号和一次负号。
\end{proof}

例如在 \(\R^2\) 上取
\[
 X=\frac{\partial}{\partial x},\qquad
 Y=x\frac{\partial}{\partial y}.
\]
则 \(X\) 沿 \(x\) 方向平移，而 \(Y\) 以依赖于 \(x\) 的速度沿 \(y\) 方向
移动；坐标公式给
\[
 [X,Y]=\frac{\partial}{\partial y}.
\]
先沿 \(X\) 再沿 \(Y\) 与颠倒次序所得位移的一阶部分相同，差异首先出现在
混合的二阶小量中，除以两个时间参数后留下 Lie 括号。因而括号为零表达的是
两个局部运动在最低非平凡阶可交换，而非仅仅两个向量在每一点线性无关或相关。

若 \(F:M\to N\) 是微分同胚，向量场可以自然地推到 \(N\) 上：
\[
 (F_*X)_{F(p)}=dF_pX_p.
\]
一般光滑映射可能把多个点映到同一点，而不同原像处推出的向量未必相同，
因此不能无条件定义整个向量场的推前。

\section{余切空间和张量场}

\(T_pM\) 的对偶空间称为\term{余切空间}，记为 \(T_p^*M\)。若
\(f\in C^\infty(M)\)，定义
\[
 df_p(v)=v(f),\qquad v\in T_pM.
\]
于是 \(df_p\in T_p^*M\)。坐标余切向量
\((dx^1|_p,\ldots,dx^n|_p)\) 与
\((\partial/\partial x^1|_p,\ldots,\partial/\partial x^n|_p)\)
互为对偶基，因为
\[
 dx^i\left(\frac{\partial}{\partial x^j}\right)=\delta^i_j.
\]
函数的链式法则现在可以写成
\[
 d(g\circ F)_p=(dF_p)^*(dg_{F(p)}),
\]
其中 \((dF_p)^*\) 是线性映射 \(dF_p\) 的对偶映射。

切向量与余切向量虽有相同维数，却没有自然的相等关系。给定坐标可以把二者
都写成 \(n\) 个数，但换坐标时它们按互逆矩阵变换；任意线性同构
\(T_pM\to T_p^*M\) 都需要额外选择。Riemann 度量以后会用
\(v\mapsto g_p(v,\mathord\cdot)\) 给出这样的同构，没有度量时不应把上、下
指标任意互换。

令 $V=T_pM$。按照本书的约定，一个 $(r,s)$ 型张量是多重线性映射
\[
 T:(V^*)^r\times V^s\longrightarrow\R.
\]
当 $r=s=0$ 时，张量就是实数。给定 $v_1,\ldots,v_r\in V$ 和
$\alpha_1,\ldots,\alpha_s\in V^*$，记
\[
 (v_1\otimes\cdots\otimes v_r\otimes\alpha_1\otimes\cdots\otimes\alpha_s)
 (\xi_1,\ldots,\xi_r;w_1,\ldots,w_s)
 =\prod_{a=1}^r\xi_a(v_a)\prod_{b=1}^s\alpha_b(w_b).
\]
这定义了一个张量。若 $(e_i)$ 是 $V$ 的基、$(\varepsilon^i)$ 是对偶基，
把每个输入按基展开，多重线性性给
\[
 T=\sum_{i_1,\ldots,i_r,j_1,\ldots,j_s}
 T^{i_1\cdots i_r}{}_{j_1\cdots j_s}\,
 e_{i_1}\otimes\cdots\otimes e_{i_r}
 \otimes\varepsilon^{j_1}\otimes\cdots\otimes\varepsilon^{j_s},
\]
其中系数等于 $T(\varepsilon^{i_1},\ldots,\varepsilon^{i_r};
e_{j_1},\ldots,e_{j_s})$。代入这些基输入还证明展开唯一。因此这些简单
张量组成一组基，张量空间维数为 $n^{r+s}$（标量情形维数为一）。有限维
张量积记号 $V^{\otimes r}\otimes(V^*)^{\otimes s}$ 表示的正是这个空间。
若张量在各坐标基下的分量都是光滑函数，就得到一个\term{张量场}。向量场
是 \((1,0)\) 型张量场，一次形式是 \((0,1)\) 型张量场；
第\ref{chap:riemannian-geometry}章中的
Riemann 度量是对称正定的 \((0,2)\) 型张量场。

张量积使类型相加；一个上指标与一个下指标可以收缩。换坐标时，上指标和
下指标按互逆的矩阵规律变换，所以完全收缩所得的标量与坐标无关。张量本身
由多重线性映射定义，坐标指标只是其计算表示。例如收缩第一个上、下指标，
就是对 $T$ 的相应两个输入代入 $(\varepsilon^i,e_i)$ 后对 $i$ 求和。
若换基矩阵为 $A$，对偶基的变换矩阵为 $A^{-1}$，求和中的两矩阵相乘为
恒等矩阵，因此结果与所选基无关。

\section{微分形式与楔积}

\(p\) 处的交错 \(k\) 次余切张量组成空间
\(\Lambda^kT_p^*M\)。一个光滑 \(k\)-形式是在每一点选取
\(\omega_p\in\Lambda^kT_p^*M\)，并要求它在任一局部坐标中写成
\[
 \omega=\sum_{i_1<\cdots<i_k}
 a_{i_1\cdots i_k}\,
 dx^{i_1}\wedge\cdots\wedge dx^{i_k}
\]
时，各系数都光滑。全体光滑 \(k\)-形式记为 \(\Omega^k(M)\)，并约定
\(\Omega^0(M)=C^\infty(M)\)。当 \(k>n\) 时，\(\Omega^k(M)=\{0\}\)。

楔积逐点由外代数定义。若
\(\alpha\in\Omega^k(M)\)、\(\beta\in\Omega^\ell(M)\)，则
\[
 \alpha\wedge\beta=(-1)^{k\ell}\beta\wedge\alpha.
\]
若 \(F:M\to N\) 光滑，定义形式的\term{拉回}
\[
 (F^*\omega)_p(v_1,\ldots,v_k)
 =\omega_{F(p)}(dF_pv_1,\ldots,dF_pv_k).
\]
由定义可得
\[
 F^*(\alpha\wedge\beta)=F^*\alpha\wedge F^*\beta,
 \qquad (G\circ F)^*=F^*\circ G^*.
\]

若 \(\iota:S\hookrightarrow M\) 是嵌入子流形，则
\(\iota^*\omega\) 称为 \(\omega\) 在 \(S\) 上的限制。这里不仅限制
坐标系数，还要先用 \(d\iota\) 把 \(S\) 的切向量送入 \(M\) 的切空间。

拉回也包含熟悉的曲线积分换元。若
\(\alpha=P\,dx+Q\,dy\) 是 \(\R^2\) 上的一次形式，而
\(\gamma(t)=(x(t),y(t))\)，则
\[
 \gamma^*\alpha
 =\bigl(P(\gamma(t))x'(t)+Q(\gamma(t))y'(t)\bigr)\,dt.
\]
因此一次形式沿曲线的积分就是其拉回在参数区间上的积分。参数改变之所以不改
变积分值，是在保向重参数化下使用拉回的复合律和一元变量代换公式；反向
重参数化使积分变号。

\section{拉回、外微分与 Cartan 公式}

在坐标图 \((U,x)\) 中，按照欧氏空间中的公式定义
\[
 d\left(\sum_Ia_I\,dx^I\right)
 =\sum_{I,j}\partial_ja_I\,dx^j\wedge dx^I.
\]
\chapterref{chap:differential-forms-euclidean}已经证明外微分与欧氏空间中
光滑映射的拉回可交换。坐标变换正是这样的映射，因此两张图给出的定义在
重叠处一致，局部算子可以拼成全局算子
\[
 d:\Omega^k(M)\longrightarrow\Omega^{k+1}(M).
\]

在低次数中，这一公式就是熟悉的偏导组合。对函数 \(f\)，
\[
 df=\sum_i\partial_i f\,dx^i;
\]
在平面上若 \(\alpha=P\,dx+Q\,dy\)，则
\[
 d\alpha=(\partial_xQ-\partial_yP)\,dx\wedge dy.
\]
系数中的交叉偏导之差来自楔积的反对称性。下一章把 \(d\alpha\) 的积分与
\(\alpha\) 在边界上的积分联系起来。

\begin{theorem}[流形上的外微分]
\label{thm:manifold-exterior-derivative}
外微分是唯一具有下列性质的实线性局部算子：在函数上，
\[
 df(v)=v(f);
\]
若 \(\alpha\in\Omega^k(M)\)，则
\[
 d(\alpha\wedge\beta)
 =d\alpha\wedge\beta+(-1)^k\alpha\wedge d\beta;
\]
并且 \(d^2=0\)。此外，对任意光滑映射 \(F:M\to N\)，都有
\[
 d(F^*\omega)=F^*(d\omega).
\]
\end{theorem}

\begin{proof}
存在性及各恒等式都可在坐标图中检验，归结为欧氏空间中的相应结论。为证
唯一性，注意每个形式局部都是
\[
 a_I\,dx^{i_1}\wedge\cdots\wedge dx^{i_k}
\]
的有限和。算子在函数上的作用、分次 Leibniz 规则以及
\(d(dx^i)=d^2x^i=0\) 已经唯一决定每一项的外微分，因而唯一决定全局算子。
\end{proof}

对向量场 \(X\)，内乘算子
\(\iota_X:\Omega^k(M)\to\Omega^{k-1}(M)\) 在 $k\ge1$ 时定义为
\[
 (\iota_X\omega)(X_1,\ldots,X_{k-1})
 =\omega(X,X_1,\ldots,X_{k-1}).
\]
约定 $\Omega^{-1}(M)=\{0\}$、$\iota_Xf=0$。Lie 导数可以先用向量场及其括号定义：
\[
 (\mathcal L_X\omega)(X_1,\ldots,X_k)
 =X\bigl(\omega(X_1,\ldots,X_k)\bigr)
 -\sum_{j=1}^k
 \omega(X_1,\ldots,[X,X_j],\ldots,X_k),
\]
并约定 \(\mathcal L_Xf=Xf\)。这个表达确实是微分形式：将任一输入 $Y$
换为 $aY$ 时，第一项中的 $X(a)$ 项，与
$[X,aY]=a[X,Y]+X(a)Y$ 在对应括号项中产生的项相消，余下表达乘以 $a$。
因此它对每个输入都为 $C^\infty(M)$ 线性。交换两个输入使整式变号，坐标
表达中的系数光滑，故得到光滑交错张量。

\begin{theorem}[Cartan 公式]
\label{thm:cartan-magic-formula}
对任意向量场 \(X\) 和微分形式 \(\omega\)，有
\[
 \mathcal L_X\omega=d(\iota_X\omega)+\iota_X(d\omega).
\]
\end{theorem}

\begin{proof}
先核对所需的导子性质。把楔积代入 Lie 导数的定义，并分别让 \(X\) 作用于
两个因子的系数，可得
\[
 \mathcal L_X(\alpha\wedge\beta)
 =(\mathcal L_X\alpha)\wedge\beta
  +\alpha\wedge\mathcal L_X\beta .
\]
因此 \(\mathcal L_X\) 是零次导子。若 \(\alpha\in\Omega^p(M)\)，则内乘的
分次 Leibniz 规则为
\[
 \iota_X(\alpha\wedge\beta)
 =\iota_X\alpha\wedge\beta+(-1)^p\alpha\wedge\iota_X\beta.
\]
把它和外微分的分次 Leibniz 规则依次用于
\(D=d\iota_X+\iota_Xd\)，含
\(d\alpha\wedge\iota_X\beta\) 与
\(\iota_X\alpha\wedge d\beta\) 的交叉项分别抵消，得到
\[
 D(\alpha\wedge\beta)=D\alpha\wedge\beta+\alpha\wedge D\beta.
\]
所以 \(D\) 也是零次导子。局部坐标下的微分形式由光滑函数和
\(dx^i=d(x^i)\) 生成，只需在函数 \(f\) 和一次形式 \(df\) 上检验等式。

在函数上，\(\iota_Xf=0\)，所以右端为
\[
 \iota_Xdf=df(X)=Xf=\mathcal L_Xf.
\]
对 \(df\)，右端为
\[
 d(\iota_Xdf)+\iota_Xd^2f=d(Xf).
\]
另一方面，对任意向量场 \(Y\)，
\[
 (\mathcal L_Xdf)(Y)
 =X(Yf)-df([X,Y])
 =X(Yf)-[X,Y]f
 =Y(Xf),
\]
也等于 \(d(Xf)(Y)\)。两算子在局部生成元上一致，故在所有微分形式上
一致。
\end{proof}

若 \(X\) 有对时间和初值联合光滑的局部流 \(\Phi_t\)，则上述定义还满足
\[
 \mathcal L_X\omega
 =\left.\frac{d}{dt}\right|_{t=0}\Phi_t^*\omega.
\]
这给出了 Lie 导数的动力学解释。证明这一等式需要局部流对初值与参数的
光滑依赖，计算见习题\ref{ex:forms-flow-lie-derivative}及其解答；Cartan
公式本身则可由上面的代数论证直接得到。

\section*{习题}
\addcontentsline{toc}{section}{习题}

\subsection*{A组　切空间和切映射}
\begin{enumerate}[label=\textbf{\thechapter.\arabic*.}]
 \exerciseitem{tangent-coordinate-basis-change}
 从导子定义推导两组坐标之间切向量基和分量的变换公式。
 \exerciseitem{tangent-hadamard-derivation}
 补全\theoremref{thm:tangent-curves-derivations}中的 Hadamard 分解，并
 证明每个导子由坐标函数上的值唯一决定。
 \exerciseitem{tangent-map-curve-chain-rule}
 分别用曲线表示和导子表示证明切映射的链式法则。
 \exerciseitem{tangent-regular-level-kernel}
 设 \(F:M^n\to N^m\) 光滑，\(m\le n\)，而 \(c\in N\) 是正则值。若
 \(S=F^{-1}(c)\)，证明对每个 \(p\in S\)，
 \[
  T_pS=\ker dF_p .
 \]
 \exerciseitem{tangent-product-splitting}
 构造自然同构
 \[
  T_{(p,q)}(M\times N)\cong T_pM\oplus T_qN,
 \]
 并证明它与积映射的切映射相容。
\end{enumerate}

\subsection*{B组　向量场与张量}
\begin{enumerate}[label=\textbf{\thechapter.\arabic*.},resume]
 \exerciseitem{lie-bracket-coordinate-formula}
 逐项推导 Lie 括号的坐标公式，并核对二阶偏导项为何相消。
 \exerciseitem{lie-bracket-jacobi-expansion}
 不引用命题~\ref{prop:lie-bracket-properties}，把向量场当作算子展开并
 证明 Jacobi 恒等式。
 \exerciseitem{lie-bracket-diffeomorphism-naturality}
 若 \(F\) 是微分同胚，证明
 \[
  F_*[X,Y]=[F_*X,F_*Y].
 \]
 \exerciseitem{cotangent-chain-rule}
 由定义证明 \(d(g\circ F)=F^*(dg)\)，并写出它的矩阵形式。
 \exerciseitem{tensor-contraction-invariance}
 对一个 \((1,1)\) 型张量证明其收缩在换坐标后不变，并说明该收缩就是相应
 线性算子的迹。
\end{enumerate}

\subsection*{C组　微分形式与 Cartan 公式}
\begin{enumerate}[label=\textbf{\thechapter.\arabic*.},resume]
 \exerciseitem{forms-pullback-wedge-functorial}
 从定义证明拉回保持楔积并反转复合次序。
 \exerciseitem{forms-exterior-coordinate-free-formula}
 证明对 \(\omega\in\Omega^k(M)\)，
 \[
 \begin{aligned}
 d\omega(X_0,\ldots,X_k)
 ={}&\sum_i(-1)^iX_i
 \omega(X_0,\widehat X_i,\ldots,X_k)\\
 &+\sum_{i<j}(-1)^{i+j}
 \omega\bigl([X_i,X_j],X_0,\ldots,\widehat X_i,\ldots,\\
 &\hspace{10em}\widehat X_j,\ldots,X_k\bigr).
 \end{aligned}
 \]
 第二项中，省略的向量保持原来的相对次序。
 \exerciseitem{forms-cartan-one-form-direct}
 对一般一次形式直接展开并验证 Cartan 公式，不使用“由生成元推广”的论证。
 \optionalexerciseitem{forms-flow-lie-derivative}
 设光滑向量场 \(X\) 已有对时间和初值联合光滑的局部流 \(\Phi_t\)。
 在共同的小坐标邻域及时间区间内，用坐标链式法则证明
 \[
  \left.\frac{d}{dt}\right|_{t=0}\Phi_t^*f=Xf,\qquad
  \left.\frac{d}{dt}\right|_{t=0}\Phi_t^*(df)=d(Xf),
 \]
 再把结论推广到所有微分形式。
\end{enumerate}
